% Member 3 - final report, III.E Structural Mechanism Analysis
\subsection{Structural Mechanism Analysis}
\label{sec:mechanism}

Before asking whether the spatial modulation helps the classifier, we check
whether it does what it is designed to do at the image level. The two
questions are separate, and this section answers only the first.

\paragraph{The comparison.}
Three augmentation sets are generated from the same 1{,}349 content images,
with the same ten style references, the same pairing, the same network and
the same file format. They differ only in the rule that computes the
blending weight $\alpha$ of Equation~\eqref{eq:blend}:

\begin{itemize}
\item \texttt{ours}: the spatial field of Equation~\eqref{eq:alphamap}, with
$\alpha_0 = 0.5$ and $\beta = 0.7$;
\item \texttt{adain}: the constant $\alpha = 0.5$, i.e. plain AdaIN;
\item \texttt{matched}: the constant $\alpha = 0.374$, chosen so that its
average strength equals that of \texttt{ours}.
\end{itemize}

The third set is what makes the comparison meaningful. The spatial field
reduces the average strength, because it protects the strong-structure
regions; comparing \texttt{ours} against \texttt{adain} alone would conflate
the effect of the spatial distribution with the effect of the lower mean
strength. Against \texttt{matched}, the two sets differ \emph{only} in how
the strength is distributed over the image.

\paragraph{The measured mean strength.}
The average strength of the \texttt{ours} field is
$\bar{\alpha} = 0.3736 \pm 0.0273$ over the 1{,}349 images, measured from
the generated $\alpha$ maps rather than assumed from $\alpha_0$. The
nominal value $\alpha_0(1-\beta\bar{M})$ predicts $0.374$, in agreement
with the measurement. The \texttt{matched} control uses the constant
$0.374$, so the two sets have equal mean strength by construction rather
than by coincidence.

\paragraph{The EdgeScore.}
The EdgeScore is computed for all three sets over all 1{,}349 images with a
fixed Canny detector (Section~\ref{sec:structure}). The spatial field raises
the mean EdgeScore from $0.3211$ for plain AdaIN to $0.3436$ for
\texttt{ours}: it does preserve more of the source edge structure than
uniform AdaIN at the same nominal strength. The improvement is consistent
over the whole set rather than being driven by a few images, and it holds
in every one of the seven classes (Table~\ref{tab:edgeclass}).

\paragraph{The matched control is the informative one.}
The \texttt{matched} control reaches an EdgeScore of $0.3467$, which is
\emph{slightly higher} than the spatial field's $0.3436$. This is not a
contradiction but a direct consequence of the design: uniform AdaIN at a
lower strength keeps more of the original image than a modulated field
whose mean is the same, because the modulated field still applies the full
strength $\alpha_0$ in its flat regions. Preserving structure and reducing
the average style are related but distinct effects, and the matched control
separates them.

\paragraph{The mechanism test is null.}
If the modulation preserved structure \emph{for its stated reason}, the
images with the most structure should benefit the most. We test this
directly, by splitting the 1{,}349 images at the median structural density
and comparing the EdgeScore gain of \texttt{ours} over \texttt{matched} in
the two halves. It does not hold. The high-density half gains
$-0.0029 \pm 0.0084$ and the low-density half $-0.0034 \pm 0.0138$; the
Welch test between them gives $t = +0.84$, $p = 0.40$. The correlation
between an image's structural density and its gain is negligible
(Pearson $r = +0.037$, $p = 0.18$; Spearman $\rho = +0.078$,
$p = 0.004$). The spatial field does \emph{not} buy extra edge retention
where the structure is strongest.

The reason is visible in the per-class breakdown
(Table~\ref{tab:edgeclass}): the \texttt{ours} minus \texttt{matched}
difference is negative and remarkably uniform across the seven classes
($-0.0010$ for \texttt{giraffe} to $-0.0048$ for \texttt{person}),
including the classes with the highest and the lowest edge density. The
modulation therefore does not act selectively on structure. It reduces the
strength almost everywhere, and the accompanying loss of edge agreement is
roughly constant per image. We report this because it changes the
interpretation of the method: the field is not the selective structure
protector it was designed to be.

\paragraph{What this establishes, and what it does not.}
Two findings, one negative and one positive. Negatively, the spatial
modulation does not preserve structure selectively, and the mechanism
designed to explain a gain does not explain one. Positively, the matched
control shows that a \emph{uniform} AdaIN at a lower mean strength retains
more edge structure than the modulated field at the same mean strength. The
honest summary is that the modulation behaves like a mild, nearly uniform
reduction of the style strength, and that the part of its behaviour which
is not that --- the spatial modulation itself --- contributes nothing that
we can detect, either in the image or in the classifier
(Section~\ref{sec:strength_energy}).
